\begin{exo}{}
	Dans un repère orthogonal, on a tracé la représentation graphique d'une fonction $f$ définie sur un intervalle $I$.
	\vspace*{2mm}
	\begin{center}
		\begin{tikzpicture}[scale=0.55, >=Stealth]
			\draw[gray!30, thin, step=1] (-7,-5) grid (6,6);
			\draw[->, thick] (-7,0) -- (6,0);
			\draw[->, thick] (0,-5) -- (0,6);
			\draw (0,0) node[below left] {O};
			\draw[black, thick] (1,0.15) -- (1,-0.15) node[below] {I};
			\draw[black, thick] (0.15,1) -- (-0.15,1) node[left] {J};
			\foreach \x in {-6,-5,-4,-3,-2,-1,2,3,4,5} { \draw (\x,0.1) -- (\x,-0.1); }
			\foreach \y in {-4,-3,-2,-1,2,3,4} { \draw (0.1,\y) -- (-0.1,\y); }
			\draw[black, very thick] plot[smooth, tension=0.5] coordinates { (-6, 3) (-5, 0) (-4, -1) (-3, 0) (-2, 2) (-1, 4) (0, 5) (1, 4) (2, 0) (3, -4) (4, -2) (5, -1) } node[below right, xshift=-2mm] {$\Cf$};
			\fill[black] (-6, 3) circle (3pt);
			\fill[black] (-5, 0) circle (3pt);
			\fill[black] (-4, -1) circle (3pt);
			\fill[black] (-3, 0) circle (3pt);
			\fill[black] (-2, 2) circle (3pt);
			\fill[black] (-1, 4) circle (3pt);
			\fill[black] (0, 5) circle (3pt);
			\fill[black] (1, 4) circle (3pt);
			\fill[black] (2, 0) circle (3pt);
			\fill[black] (3, -4) circle (3pt);
			\fill[black] (4, -2) circle (3pt);
			\fill[black] (5, -1) circle (3pt);
		\end{tikzpicture}
	\end{center}
	\begin{enumerate}
		\item Résoudre, sur $I$, les inéquations :
		      \vspace*{-2mm}
		      \begin{multicols}{2}
		      \begin{enumerate}
			      \item $f(x) > 4$
			      \item $0 \leqslant f(x) \leqslant 4$
		      \end{enumerate}
		      \end{multicols}
		      \vspace*{-2mm}
		\item Indiquer le signe de $f(x)$ suivant les valeurs de $x$.
	\end{enumerate}

\end{exo}
